\chapter{Methotrexate Level}
\section*{What Is This Test?}
The methotrexate level measures the serum concentration of methotrexate (MTX), an antimetabolite chemotherapeutic and immunosuppressive agent. MTX is a folic acid antagonist that competitively inhibits dihydrofolate reductase (DHFR), depleting tetrahydrofolate (THF) cofactors, thereby blocking de novo purine and thymidylate synthesis and inhibiting DNA synthesis and cell division \cite{bleyer1978clinical}. At low doses (oral or subcutaneous, 7.5--25 mg weekly), MTX is used for autoimmune diseases (rheumatoid arthritis, psoriasis, psoriatic arthritis, inflammatory bowel disease, systemic lupus erythematosus). At high doses (IV, 1--12 g/m$^2$), MTX is used as a chemotherapeutic agent for certain cancers (acute lymphoblastic leukemia/ALL, non-Hodgkin lymphoma, primary CNS lymphoma, osteosarcoma, choriocarcinoma, breast cancer, and head and neck cancers). High-dose MTX requires rigorous TDM and leucovorin (folinic acid) rescue to prevent life-threatening toxicity \cite{jolivet1983pharmacology}. MTX is cleared primarily by renal excretion (90\% unchanged in urine). Delayed MTX clearance (due to renal impairment, third-space fluid collections [ascites, pleural effusions], drug interactions [NSAIDs, penicillins, probenecid, proton pump inhibitors], or acidic urine pH) produces prolonged supratherapeutic levels that cause severe toxicity: myelosuppression (pancreatopenia with nadir 7--14 days), mucositis (the most common and earliest sign of severe MTX toxicity --- severe oral and GI mucositis, hemorrhagic enteritis), acute kidney injury (intratubular crystallization of MTX and its metabolite 7-hydroxy-MTX), hepatitis (transient transaminase elevation with high-dose MTX), and neurotoxicity (acute encephalopathy, leukoencephalopathy with chronic administration) \cite{howard2016preventing}.
\section*{Alternative Names}
MTX Level; Methotrexate Serum Concentration; MTX Drug Level; High-Dose Methotrexate Level
\section*{Principle}
Methotrexate is measured by particle-enhanced turbidimetric inhibition immunoassay (PETINIA) or fluorescence polarization immunoassay (FPIA) on automated analyzers. LC-MS/MS is the gold standard and is preferred because immunoassays cross-react with MTX metabolites (7-hydroxy-MTX and DAMPA --- 2,4-diamino-N10-methylpteroic acid) that accumulate in renal impairment, overestimating MTX concentrations. Carboxypeptidase G2 (glucarpidase, Voraxaze), the MTX detoxifying enzyme used as rescue for severe MTX toxicity, rapidly cleaves MTX to inactive metabolites and interferes with most MTX immunoassays, producing falsely undetectable results \cite{green2012glucarpidase}.
\section*{History}
Methotrexate (then called amethopterin) was synthesized in the late 1940s by Yellapragada Subbarow and colleagues at Lederle Laboratories as an analog of aminopterin, the antifolate that Sidney Farber showed in 1948 could produce remissions in childhood acute lymphoblastic leukemia. Leucovorin (citrovorum factor) rescue, the principle that permits safe administration of high-dose methotrexate, was established experimentally by Goldin and colleagues in the 1950s. High-dose methotrexate with leucovorin rescue entered clinical oncology in the 1970s for osteosarcoma and hematologic malignancies. Low-dose weekly methotrexate was approved for psoriasis in the early 1970s and for rheumatoid arthritis in 1988, becoming the anchor drug for the treatment of rheumatoid arthritis. Therapeutic drug monitoring of high-dose methotrexate has been standard practice since the 1970s.
\section*{Why This Test Is Ordered}
Mandatory TDM after high-dose MTX administration to determine when it is safe to discontinue leucovorin rescue. Serum MTX levels are measured at 24, 48, and 72 hours after the start of the MTX infusion. The level at 48 hours is the primary determinant of the required leucovorin dose and duration. Leucovorin rescue is continued until MTX is <0.05--0.1 mcmol/L. TDM is not routinely performed for low-dose weekly MTX in autoimmune disease.
\section*{Primary vs Secondary Test}
This is a secondary, therapeutic-monitoring test rather than a diagnostic test. Methotrexate levels are not diagnostic; they are measured during and after high-dose methotrexate infusions to identify delayed clearance and to determine when leucovorin rescue can be safely discontinued.
\section*{How This Test Is Performed}
Blood is collected by standard venipuncture into a plain serum tube (red-top) or a serum separator tube (SST, gold-top). Samples are drawn at scheduled time points measured from the start of the infusion, most commonly at 24, 48, and 72 hours. The serum methotrexate concentration is measured by particle-enhanced turbidimetric inhibition immunoassay (PETINIA), fluorescence polarization immunoassay (FPIA), or LC-MS/MS on automated analyzers; LC-MS/MS is preferred because immunoassays cross-react with the metabolites 7-hydroxy-methotrexate and DAMPA that accumulate in renal impairment. Results are reported in mcmol/L and are typically available within 1--2 hours because dosing decisions depend on them.
\section*{What Preparation Is Needed}
No fasting or special preparation is required. The critical requirement is precise timing: levels are drawn at fixed hours (24, 48, 72) after the start of the infusion, so the infusion start time must be documented accurately. Hydration and urinary alkalinization are maintained per protocol (urine pH above 7.0), and the dose and infusion schedule must be consistent with the leucovorin rescue plan.
\section*{Contraindications and Precautions}
\begin{itemize}
  \item Recognize the early signs of methotrexate toxicity: mucositis (earliest sign), myelosuppression, renal impairment, and hepatitis
  \item Identify patients at high risk of delayed clearance: renal impairment, third-space fluid collections (ascites, pleural effusions), acidic urine, and dehydration
  \item Avoid drugs that delay methotrexate clearance: NSAIDs, penicillins, probenecid, and proton pump inhibitors
  \item Interpret immunoassay results cautiously after glucarpidase (carboxypeptidase G2) administration, which rapidly cleaves methotrexate and causes falsely low immunoassay results; collect samples at least 24 hours after glucarpidase and confirm by LC-MS/MS if needed
  \item Continue leucovorin rescue until the methotrexate level falls below 0.05--0.1 mcmol/L; do not discontinue rescue based on clinical appearance alone
\end{itemize}
\section*{Risks}
Minimal risk. Slight pain or bruising at the venipuncture site. Rarely: hematoma, infection, or vasovagal syncope.
\section*{What Happens After the Test}
Each methotrexate level is compared with the expected clearance curve and the time-specific thresholds (for example, below 1 mcmol/L at 48 hours). A level above the threshold for the time point triggers continued or intensified leucovorin rescue, aggressive IV hydration, and urinary alkalinization, with a repeat level within 24 hours. If a high level is accompanied by renal dysfunction, glucarpidase (Voraxaze) may be given \cite{widemann2006understanding}. Levels are repeated until methotrexate is below 0.05--0.1 mcmol/L, at which point leucovorin rescue is discontinued.
\section*{Understanding Results}
Elevated 24-hour MTX: >5--10 mcmol/L. Elevated 48-hour MTX: >1 mcmol/L (increased risk of myelosuppression and mucositis); >5 mcmol/L (severe, prolonged myelosuppression). Elevated 72-hour MTX: >0.1 mcmol/L (requires continued leucovorin rescue). Critical delayed elimination at any time point requires: aggressive IV hydration (3 L/m$^2$/day) with urinary alkalinization (sodium bicarbonate to maintain urine pH >7.0), increased leucovorin dose and frequency, and consideration of glucarpidase (Voraxaze) if MTX is >30 mcmol/L at 24 hours, >10 mcmol/L at 48 hours, or >5 mcmol/L at 72 hours AND renal dysfunction is present or developing.
\section*{Normal Results}
Expected MTX levels after high-dose MTX (timed from start of infusion): at 24 hours: <5--10 mcmol/L; at 48 hours: <0.5--1.0 mcmol/L; at 72 hours: <0.05--0.1 mcmol/L (goal for discontinuation of leucovorin rescue). Low-dose MTX (autoimmune disease): TDM not routinely performed.
\section*{Follow-up Tests}
Serum creatinine and BUN (before each high-dose MTX dose and daily during clearance monitoring), urine pH (maintain >7.0), complete blood count (at nadir, 7--14 days post-dose), liver function tests, and adequate hydration/urinary output monitoring.
\section*{References}
\bibliographystyle{unsrtnat}
\bibliography{references}

\clearpage
